\section*{Розділ III. Процедури проведення засідань}
\addcontentsline{toc}{section}{Розділ III. Процедури проведення засідань}

\subsection*{3.1. Відкриття та ведення засідання}
\addcontentsline{toc}{subsection}{3.1. Відкриття та ведення засідання}
    3.1.1. Засідання КСУ відкривається Спікером КСУ у визначений час після проведення реєстрації учасників.

    3.1.2. Процедура відкриття засідання включає такі етапи:

        \begin{enumerate}[label=\alph*)]
            \item Реєстрація учасників засідання та технічна перевірка засобів зв'язку (у разі дистанційної участі);
            \item Встановлення кворуму та констатація правомочності засідання;
            \item Представлення запрошених осіб (представників адміністрації Університету, експертів, консультантів) та оголошення статусу їх участі;
            \item Регламентні повідомлення щодо технічних особливостей проведення засідання, процедурних правил, змін у складі делегатів та інформації про стан виконання попередніх рішень КСУ;
            \item Затвердження порядку денного засідання.
        \end{enumerate}

    3.1.3. Засідання КСУ веде Спікер КСУ. У разі неможливості виконання Спікером своїх обов'язків, засідання може вести особа, уповноважена КСУ.

\subsection*{3.2. Кворум та правомочність КСУ}
\addcontentsline{toc}{subsection}{3.2. Кворум та правомочність КСУ}
    3.2.1. Засідання КСУ є правомочним (має кворум), якщо на ньому присутні не менше половини від загального складу делегатів КСУ відповідно до пункту 3.4.1 Положення про ОСС.

    3.2.2. Кворум перевіряється на початку засідання та може бути перевірений повторно на вимогу будь-якого делегата КСУ.

    3.2.3. У разі втрати кворуму під час засідання, Спікер КСУ зобов'язаний припинити розгляд питань до відновлення кворуму або закрити засідання.

    3.2.4. При проведенні засідання у дистанційному форматі, присутність делегата підтверджується його активною участю у засіданні (підключення до відеоконференції, участь у голосуванні тощо).

\subsection*{3.3. Порядок денний засідання}
\addcontentsline{toc}{subsection}{3.3. Порядок денний засідання}
    3.3.1. Проєкт порядку денного, підготовлений Спікером КСУ, виноситься на затвердження КСУ на початку засідання.

    3.3.2. КСУ може внести зміни до проєкту порядку денного простою більшістю голосів від присутніх делегатів, включаючи:

        \begin{enumerate}[label=\alph*)]
            \item Додавання нових питань до порядку денного;
            \item Виключення питань з порядку денного;
            \item Зміну послідовності розгляду питань.
        \end{enumerate}

    3.3.3. Пропозиції щодо змін до порядку денного можуть вноситися делегатами КСУ до його затвердження. Кожна пропозиція обговорюється та виноситься на голосування окремо.

    3.3.4. Після затвердження порядку денного його зміна можлива лише за рішенням КСУ, прийнятим простою більшістю голосів від присутніх делегатів.

\subsection*{3.4. Процедури обговорення питань}
\addcontentsline{toc}{subsection}{3.4. Процедури обговорення питань}
    3.4.1. Розгляд кожного питання порядку денного включає:

        \begin{enumerate}[label=\alph*)]
            \item Доповідь з питання;
            \item Відповіді на запитання до доповідача;
            \item Обговорення (дискусію);
            \item Прийняття рішення.
        \end{enumerate}

    3.4.2. Доповідь з питання робить особа, визначена порядком денним, або за дорученням Спікера КСУ. Якщо доповідач не визначений, доповідь робить ініціатор розгляду питання.

    3.4.3. Після доповіді делегати КСУ можуть ставити запитання доповідачу. Запитання мають стосуватися суті доповіді та бути сформульованими стисло.

    3.4.4. Обговорення питання проводиться у формі дискусії, під час якої делегати КСУ можуть висловлювати свої позиції, пропозиції та зауваження з питання.

\subsection*{3.5. Права та обов'язки делегатів}
\addcontentsline{toc}{subsection}{3.5. Права та обов'язки делегатів}
    3.5.1. Кожен делегат КСУ має право:

        \begin{enumerate}[label=\alph*)]
            \item Брати участь в обговоренні всіх питань порядку денного;
            \item Вносити пропозиції та поправки до проєктів рішень;
            \item Ставити запитання доповідачам та іншим учасникам засідання;
            \item Голосувати з усіх питань, що виносяться на голосування;
            \item Вимагати перевірки кворуму;
            \item Звертатися до Спікера КСУ з процедурних питань.
        \end{enumerate}

    3.5.2. Кожен делегат КСУ зобов'язаний:

        \begin{enumerate}[label=\alph*)]
            \item Дотримуватися регламенту засідання;
            \item Не перешкоджати роботі засідання;
            \item Поводитися коректно щодо інших учасників засідання;
            \item Вимикати мікрофон, коли не виступає (при дистанційній участі).
        \end{enumerate}

\subsection*{3.6. Регламент виступів}
\addcontentsline{toc}{subsection}{3.6. Регламент виступів}
    3.6.1. Для виступів на засіданні КСУ встановлюється такий регламент:

        \begin{enumerate}[label=\alph*)]
            \item Доповідь з питання порядку денного -- до 15 хвилин;
            \item Виступ у дискусії -- до 5 хвилин;
            \item Відповідь на запитання -- до 2 хвилин;
            \item Репліка -- до 1 хвилини.
        \end{enumerate}

    3.6.2. Спікер КСУ може продовжити час виступу за мотивованим проханням виступаючого або за рішенням КСУ, прийнятим простою більшістю голосів від присутніх делегатів.

    3.6.3. Спікер КСУ попереджає виступаючого про закінчення часу за 1 хвилину до його завершення та має право зупинити виступ після закінчення регламентного часу.

    3.6.4. Порядок надання слова визначається Спікером КСУ з урахуванням:

        \begin{enumerate}[label=\alph*)]
            \item Черговості подання заявок на виступ;
            \item Необхідності забезпечення збалансованої дискусії;
            \item Представництва різних точок зору з питання.
        \end{enumerate}

\subsection*{3.7. Процедури підтримання порядку}
\addcontentsline{toc}{subsection}{3.7. Процедури підтримання порядку}
    3.7.1. За порушення регламенту засідання або неналежну поведінку Спікер КСУ може застосувати такі заходи:

        \begin{enumerate}[label=\alph*)]
            \item Попередження делегата про порушення;
            \item Позбавлення слова;
            \item Видалення делегата із засідання за систематичні порушення порядку.
        \end{enumerate}

    3.7.2. При проведенні засідання у дистанційному форматі Спікер КСУ має право:

        \begin{enumerate}[label=\alph*)]
            \item Відключити мікрофон учасника, який порушує порядок;
            \item Видалити учасника із відеоконференції за систематичні порушення;
            \item Припинити трансляцію засідання у разі технічних проблем або серйозних порушень порядку.
        \end{enumerate}

    3.7.3. Рішення про видалення делегата із засідання може бути оскаржене до КСУ і розглядається негайно простою більшістю голосів.

\subsection*{3.8. Завершення засідання}
\addcontentsline{toc}{subsection}{3.8. Завершення засідання}
    3.8.1. Засідання КСУ завершується після розгляду всіх питань порядку денного або за рішенням КСУ про дострокове завершення засідання.

    3.8.2. Перед завершенням засідання Спікер КСУ:

        \begin{enumerate}[label=\alph*)]
            \item Підводить підсумки засідання;
            \item Інформує про прийняті рішення;
            \item Визначає відповідальних за виконання рішень та строки їх виконання;
            \item Інформує про дату наступного засідання (за наявності відповідного рішення).
        \end{enumerate}

    3.8.3. Засідання вважається завершеним з моменту оголошення Спікером КСУ про його закриття.